\scanpage{60}
% archive foliation: 59
% lines 1--6, cancelled block
\deleted{
$C=(C_0,C_1)$ groupoïde connexe

$G=\operatorname{Aut}(C)$ groupe des automorphismes

$G^0$ \unsure{ss-gps} des $g\in G$ induisant l'identité sur $C_0$

\diagram{\begin{tikzcd}[ampersand replacement=\&,column sep=large]
1\arrow[r,dash]\&G^0\arrow[r]\&G\arrow[r,"{\uuline{\text{surj}}}"']\&\unsure{\mathfrak{S}_{C_0}|\mathcal{S}_{C_0}}\arrow[r]\&1
\end{tikzcd}}
}
% lines 7--12, cancelled bracketed paragraph
\deleted{[Pour qu'une autoéquivalence \inserted{$u$} soit un iso,
il f. et s. qu'elle soit bijective sur les objets de $C$, et induise un
iso. \deleted{\unsure{sur les}} $\operatorname{Aut}(\unsure{x|a})\longrightarrow\operatorname{Aut}(u_0x)$
pour \uline{un} $x\in C_0$ ; d'autre part, pour toute application
$C_0\overset{u_0}{\unsure{-|\longrightarrow}}C'_0$}

% line 13
\diagram{\begin{tikzcd}[ampersand replacement=\&,column sep=large]
1\arrow[r,dash]\&\operatorname{Aut}^0(C)\arrow[r,dash]\&\operatorname{Aut}C\arrow[r,dash]\&\unsure{\mathfrak{S}_{C}|\mathcal{S}_{C}}\arrow[r,dash]\&1
\end{tikzcd}}
% lines 14--16
\deleted{$\operatorname{Aut}^0(C)\simeq$}

$\unsure{a|0}-\,\cdot$
\diagram{\begin{tikzcd}[ampersand replacement=\&,column sep=large]
1\arrow[r]\&\operatorname{Aut}^0(C,a)\arrow[r,dash]\&\operatorname{Aut}^0(C)\arrow[r]\&\operatorname{Aut}(\underset{\substack{\Vert\\\operatorname{Aut}_C(a)}}{\Pi_a})\arrow[r]\&1
\end{tikzcd}}
\[\unsure{1}\]
% lines 17--19
\deleted{\unsure{À}} Pour déterminer
\deleted{$\operatorname{Aut}^0(C,a)=$}
\[
\operatorname{Aut}^0(C,a)=\{f\in\operatorname{Aut}(C)\mid
\operatorname{Ob}f=\operatorname{id}\ .\
(f_a:\operatorname{Aut}(a)-\operatorname{Aut}(\underset{\substack{\Vert\\ f(a)}}{a}))=\operatorname{id}\}
\]
% lines 20--22
On considère, pour $\forall$ $x\in C_0$,
\[
\unsure{T_x|T_a}=\operatorname{Isom}(a,x),\qquad\Pi_a\text{-torseur à droite}
\]
et son commutant \deleted{$\unsure{G_x=}$} $\operatorname{Aut}(x)$
